% Chapter: From Porous Media to Pore-Network Graphs

\section{The Porous Reactor Picture}
\label{sec:porous_picture}

Consider a macroscopic block of porous solid saturated with a reacting fluid.
The solid matrix occupies most of the volume; the remaining void space is a connected network of pores joined by narrow throats.
When a chemical excitation is injected at one face of the block, it enters the pore network as a pulse of concentration and propagates through the fluid.
At each throat the pulse travels at a speed set by the local chemistry and geometry; at each junction it may split, merge, annihilate with another pulse, or be blocked by refractory residue from a previous pulse.
The output is read as a concentration signal at probes on another face.

This is the physical picture that motivates the thesis.
The internal geometry is too fine to resolve in a full three-dimensional continuum simulation of a laboratory-scale block, and the connectivity is usually unknown.
The central move is therefore to abstract the pore space as a graph: bulbous pores become nodes and narrow throats become edges.
The fluid dynamics inside each element are replaced by measured transfer functions, and the network-level traffic is predicted by composing those functions.

The same abstraction is common in petroleum engineering and groundwater flow, where pore-network models predict multiphase flow and transport by assigning conductances to throats and solving for pressure and saturation on the graph \cite{xiong2016review, blunt2017multiphase}.
What changes here is the physics: the transported quantity is not a passive scalar carried by pressure-driven flow, but a self-sustaining excitation pulse whose existence depends on the local reaction--diffusion state.
The conductance analogue is a pulse transfer function, and the nodes carry memory in the form of refractory time.

\section{Why Two-Dimensional Controlled Simulations}
\label{sec:porous_why2d}

A real porous reactor is three-dimensional, so why study two-dimensional channels?
The answer is control.
In a thin quasi-two-dimensional layer of the photosensitive BZ reaction, a projected light mask defines the suppression field $\phi(x,y)$ everywhere in the domain.
The local kinetic parameters are therefore known and uniform across the layer, and the pulse-traffic rules can be isolated and measured.
In a three-dimensional bulk, absorption and scattering make the interior light field unknown, so the local kinetics would be uncontrolled.

Two-dimensional simulations are thus the controlled micro-scale limit of the three-dimensional porous problem.
They provide the local operators from which a three-dimensional pore-network model can later be calibrated once a volumetric illumination strategy is chosen.
The results in this thesis are therefore not a toy model; they are the micro-scale calibration layer that any predictive porous-medium description must rest on.

\section{Pore-Network Graph Abstraction}
\label{sec:porous_graph}

The graph abstraction is straightforward.
A pore is represented by a node; a throat by a directed or undirected edge.
Each edge carries a transit time, the time for a pulse to travel the throat length at the measured pulse speed.
Each node carries a refractory time, the interval after a pulse arrives during which the node will ignore further pulses.
Some nodes also carry an inhibition window: two pulses arriving within this window annihilate.
Finally, a routing table at each node decides which incoming edges feed which outgoing edges, encoding the junction topology.

This abstraction reduces a continuous PDE problem to a discrete event simulation.
It is fast enough to explore candidate pore networks in milliseconds, before committing to an expensive PDE simulation.
Its validity is the main question of the thesis.

\section{Local Operators}
\label{sec:porous_operators}

Four local operators are measured in the two-dimensional simulations reported in \refChap{chap:results}:

\begin{itemize}
    \item \textbf{Wire:} a straight throat transmits a pulse at a fixed speed with little attenuation.
    \item \textbf{Low-pass filter:} a throat drops pulses that arrive faster than the refractory time permits.
    \item \textbf{Inhibition:} at a T-junction, a pulse on one arm can block a pulse on another arm if it arrives within a short window.
    \item \textbf{Anisotropic routing:} a tilted diffusion tensor steers pulses along preferred directions, introducing path-dependent delays.
\end{itemize}

These four operators are sufficient to build a simple graph simulator.
More complex behaviours, such as one-way diodes or frequency-selective band-pass filters, are discussed in the literature but are not included as validated primitives here.
A diode pilot reported in \refSec{sec:res_diode} shows that a simple asymmetric wall geometry or static light barrier does not give robust rectification in the present parameter regime, so the diode is left as an open problem.

\section{Related Work on Pore-Network Models of Transport}
\label{sec:porous_related}

Pore-network models have been used for reactive transport, where species diffuse and react inside a network of pores \cite{xiong2016review, meakin2009modeling}.
The emphasis in that literature is on effective reaction rates and dispersion, not on excitable pulses.
Excitable media on networks have also been studied, usually by coupling discrete excitable elements according to a graph topology \cite{kinouchi2006optimal}.
The present work differs in two ways.
First, the local element is a continuous reaction--diffusion segment, not a discrete cellular-automaton node.
Second, the graph parameters are not chosen by hand but extracted from PDE simulations, so the model is calibrated rather than assumed.

The closest conceptual precedent is the idea of extracting discrete circuit rules from continuous excitable media.
Gorecki \etal designed specific junction geometries to obtain logic gates and information processing devices \cite{gorecki2009information}, and Adamatzky demonstrated collision-based computing in structured BZ media \cite{adamatzky2002collision}.
Those works build one geometry for one function.
This thesis builds a general composition framework: measure the local rules, put them in a graph simulator, and validate the composition against PDE solutions.
